% =============================================================================
% APPENDIX LO: LIT-LOEWENSTEIN: Visceral Influences and Emotion in Decision Making
% =============================================================================
% Module: BCM2_LIT_LOEWENSTEIN
% Category: LIT
% Version: 1.0 (January 2026)
% Status: Final
%
% Comprehensive integration of George Loewenstein's research contributions
% on visceral influences, emotions, curiosity, hot-cold empathy gaps,
% and intertemporal choice into the EBF framework.
%
% Language: English
% =============================================================================

% =============================================================================
% CROSS-REFERENCE STRUCTURE
% =============================================================================
%
% CHAPTER LINKAGE (Required):
% - Primary Chapter: Chapter 11: Awareness (CORE-AWARE)
% - Secondary Chapters: Chapter 9: Context, Chapter 13: Behavioral Change Journey
%
% This appendix integrates with:
%
% DEPENDENCIES (what this appendix REQUIRES):
% - Appendix ACE (LIT-KAHNEMAN): Dual-system framework
% - Appendix UNMAPPED_AU (CORE-AWARE): Awareness function
% - Appendix UNMAPPED_V (CORE-WHEN): Context framework
% - Appendix UNMAPPED_BBB (CORE-WHERE): Parameter repository
%
% DEPENDENTS (what REQUIRES this appendix):
% - Appendix UNMAPPED_BBB (CORE-WHERE): Parameters for LLMMC
% - Appendix UNMAPPED_AW (CORE-STAGE): Behavioral change dynamics
% - All DOMAIN appendices using emotional/visceral factors
%
% =============================================================================

\begin{tcolorbox}[colback=gray!5!white, colframe=gray!75!black,
    title=Appendix Metadata]
\textbf{Appendix:} LO --- LIT-LOEWENSTEIN \\
\textbf{Category:} Literature Integration \\
\textbf{Author Focus:} George Loewenstein (Carnegie Mellon University) \\
\textbf{Papers:} 29 (93\% Tier-1 Evidence) \\
\textbf{Primary Contribution:} Visceral influences, hot-cold empathy gap, curiosity, emotions in decision making
\end{tcolorbox}

\begin{tcolorbox}[colback=purple!5!white, colframe=purple!75!black,
    title=Chapter Linkage]

\textbf{Primary Chapter:} Chapter 11: Awareness
\begin{itemize}[nosep]
\item Section 11.4: Emotional Modulation $\leftrightarrow$ This Appendix Section \ref{sec:lo-visceral}
\item Section 11.7: State-Dependent Awareness $\leftrightarrow$ This Appendix Section \ref{sec:lo-hotcold}
\end{itemize}

\textbf{Secondary Chapters:}
\begin{itemize}[nosep]
\item Chapter 9: Context --- visceral states as $\Psi$ modifiers
\item Chapter 13: BCJ --- emotions in behavioral change
\end{itemize}

\textbf{Reading Guide:}
\begin{itemize}[nosep]
\item For conceptual overview: Read Chapter 11 first
\item For formal details: This appendix provides complete specification
\item For integration: See Appendix ACE (LIT-KAHNEMAN) on dual systems
\end{itemize}

\end{tcolorbox}


\chapter{LIT-LOEWENSTEIN: Visceral Influences and Emotion in Decision Making}
\label{app:lit-loewenstein}


\begin{abstract}
This appendix integrates George Loewenstein's research program into the EBF framework.
Loewenstein's work establishes that: (1) visceral states (hunger, pain, arousal, emotions)
systematically override cognitive preferences; (2) people in ``cold'' states cannot
accurately predict their behavior in ``hot'' states (hot-cold empathy gap); (3) curiosity
arises from information gaps; (4) anticipation generates utility independent of consumption.
His contributions fundamentally reshape how EBF models the Awareness function $A(\cdot)$
and Context $\Psi$. 29 papers spanning 1987-2016 are integrated, with 93\% providing
Tier-1 evidence.
\end{abstract}


\begin{tcolorbox}[colback=blue!5!white, colframe=blue!75!black,
    title=Quick Reference: LIT-LOEWENSTEIN]

\textbf{Core Question:} How do visceral states and emotions systematically bias decisions away from long-term preferences?

\textbf{Key Formula (Visceral Override):}
\begin{equation}
U_{\text{actual}}(t) = \alpha(V_t) \cdot U_{\text{cognitive}} + (1-\alpha(V_t)) \cdot U_{\text{visceral}}
\end{equation}
where $V_t$ is visceral intensity and $\alpha(V_t) \to 0$ as $V_t \to 1$.

\textbf{Key Concepts:}
\begin{itemize}[nosep]
\item \textbf{Visceral influences}: Bodily states that override cognition
\item \textbf{Hot-cold empathy gap}: Prediction failures across states
\item \textbf{Curiosity}: Motivation from information gaps
\item \textbf{Anticipatory utility}: Value from expecting future events
\end{itemize}

\textbf{Cross-References:} Appendix ACE (LIT-KAHNEMAN), AU (CORE-AWARE), V (CORE-WHEN), BBB (CORE-WHERE)
\end{tcolorbox}

% -----------------------------------------------------------------------------
% APPENDIX SCOPE BOX
% -----------------------------------------------------------------------------

\begin{tcolorbox}[
    colback=orange!5!white,
    colframe=orange!75!black,
    title={\textbf{Appendix Scope: Ziel / In-Scope / Out-of-Scope / Constraints}},
    fonttitle=\bfseries
]

\textbf{Ziel (Objective):}
\begin{itemize}[nosep]
    \item How do visceral states and emotions systematically override cognitive preferences, and why do people fail to predict this influence?
\end{itemize}

\textbf{In-Scope:}
\begin{itemize}[nosep]
    \item Theoretical Framework
    \item Axioms and Formal Definitions
    \item Key Results and Findings
    \item Theme 1: Visceral Influences on Behavior
    \item Theme 2: Hot-Cold Empathy Gap
\end{itemize}

\textbf{Out-of-Scope (delegiert an andere Appendices):}
\begin{itemize}[nosep]
    \item LIT-KAHNEMAN $\rightarrow$ Appendix ACE
    \item CORE-AWARE $\rightarrow$ Appendix UNMAPPED_AU
    \item CORE-WHEN $\rightarrow$ Appendix UNMAPPED_V
    \item CORE-WHERE $\rightarrow$ Appendix UNMAPPED_BBB
    \item CORE-WHERE $\rightarrow$ Appendix UNMAPPED_BBB
\end{itemize}

\textbf{Constraints (Anwendungsgrenzen):}
\begin{itemize}[nosep]
    \item [TODO: Define when this appendix does NOT apply]
    \item [TODO: Define required prerequisites]
    \item [TODO: Define known limitations]
\end{itemize}

\textbf{Lieferobjekte (Deliverables):}
\begin{itemize}[nosep]
    \item 6 Table(s)
    \item 18 Equation(s)
    \item 6 Axiom(s)
    \item Worked Example(s)
\end{itemize}

\end{tcolorbox}




\begin{tcolorbox}[colback=red!5!white, colframe=red!75!black,
    title=The Fundamental Question]
\textbf{How do visceral states and emotions systematically override cognitive preferences, and why do people fail to predict this influence?}
\end{tcolorbox}


%%%%%%%%%%%%%%%%%%%%%%%%%%%%%%%%%%%%%%%%%%%%%%%%%%%%%%%%%%%%%%%%%%%%%%%%%%%%%%%
\section{Theoretical Framework}
\label{sec:lo-theory}
%%%%%%%%%%%%%%%%%%%%%%%%%%%%%%%%%%%%%%%%%%%%%%%%%%%%%%%%%%%%%%%%%%%%%%%%%%%%%%%

Loewenstein's research program addresses three fundamental gaps in behavioral economics:

\subsection{The Emotion-Cognition Gap}

Standard economic models treat utility as purely cognitive. Loewenstein shows emotions are not just ``noise'' but systematic influences:
\begin{equation}
U_{\text{total}} = U_{\text{cognitive}} + U_{\text{emotional}} + \gamma_{CE} \cdot U_C \cdot U_E
\end{equation}
where $\gamma_{CE}$ captures the interaction between cognitive and emotional utility.

\subsection{The Prediction Gap}

People systematically fail to predict their own behavior across visceral states:
\begin{equation}
\mathbb{E}_{\text{cold}}[\text{Behavior}_{\text{hot}}] \neq \text{Behavior}_{\text{hot}}
\end{equation}
This ``hot-cold empathy gap'' explains failures in commitment devices, health behaviors, and addiction.

\subsection{The Motivation Gap}

Standard models struggle to explain curiosity, anticipation, and intrinsic motivation. Loewenstein provides:
\begin{equation}
\text{Curiosity} = f(\text{Information Gap}) \quad \text{(inverted U-shape)}
\end{equation}


%%%%%%%%%%%%%%%%%%%%%%%%%%%%%%%%%%%%%%%%%%%%%%%%%%%%%%%%%%%%%%%%%%%%%%%%%%%%%%%
\section{Axioms and Formal Definitions}
\label{sec:lo-axioms}
%%%%%%%%%%%%%%%%%%%%%%%%%%%%%%%%%%%%%%%%%%%%%%%%%%%%%%%%%%%%%%%%%%%%%%%%%%%%%%%

Based on Loewenstein's empirical findings, we formalize six axioms for the EBF:

\begin{tcolorbox}[colback=gray!5!white, colframe=gray!75!black,
    title=Axiom UNMAPPED_LO-1: Visceral Override]
\textbf{Statement:} As visceral intensity increases, cognitive preferences are increasingly dominated by immediate visceral drives.
\begin{equation}
\lim_{V \to 1} \frac{\partial U_{\text{actual}}}{\partial U_{\text{cognitive}}} = 0
\end{equation}
\textbf{Interpretation:} Extreme hunger, pain, or arousal can completely override long-term goals.

\textbf{Source:} \citet{loewenstein1996out}, \citet{loewenstein2005visceral}
\end{tcolorbox}

\begin{tcolorbox}[colback=gray!5!white, colframe=gray!75!black,
    title=Axiom UNMAPPED_LO-2: Hot-Cold Empathy Gap]
\textbf{Statement:} People in ``cold'' (non-aroused) states underestimate the influence of ``hot'' (aroused) states on their future behavior.
\begin{equation}
\mathbb{E}_{\text{cold}}[V_{\text{hot}}] < V_{\text{hot}} \quad \text{(systematic underestimation)}
\end{equation}
\textbf{Interpretation:} Commitment devices fail because people underpredict visceral temptation.

\textbf{Source:} \citet{loewenstein2006hot}, \citet{loewenstein2000projecting}
\end{tcolorbox}

\begin{tcolorbox}[colback=gray!5!white, colframe=gray!75!black,
    title=Axiom UNMAPPED_LO-3: Curiosity as Information Gap]
\textbf{Statement:} Curiosity follows an inverted U-shape with respect to knowledge level.
\begin{equation}
\text{Curiosity}(K) = K \cdot (1-K) \quad \text{where } K \in [0,1] \text{ is knowledge level}
\end{equation}
\textbf{Interpretation:} Maximum curiosity at intermediate knowledge; complete ignorance or complete knowledge suppresses curiosity.

\textbf{Source:} \citet{loewenstein1999becausecuriosity}
\end{tcolorbox}

\begin{tcolorbox}[colback=gray!5!white, colframe=gray!75!black,
    title=Axiom UNMAPPED_LO-4: Anticipatory Utility]
\textbf{Statement:} People derive utility from anticipating future consumption, independent of actual consumption.
\begin{equation}
U_{\text{total}}(t_0, t_1) = U_{\text{anticipation}}(t_0, t_1) + \delta^{t_1-t_0} \cdot U_{\text{consumption}}(t_1)
\end{equation}
\textbf{Interpretation:} Delay can increase total utility through savoring; explains preference for improving sequences.

\textbf{Source:} \citet{loewenstein1987emotions}
\end{tcolorbox}

\begin{tcolorbox}[colback=gray!5!white, colframe=gray!75!black,
    title=Axiom UNMAPPED_LO-5: Risk as Feelings]
\textbf{Statement:} Risk assessment is driven by emotional reactions, not just cognitive probability weighting.
\begin{equation}
\text{Risk Perception} = \omega_C \cdot \text{Cognitive Assessment} + \omega_E \cdot \text{Emotional Reaction}
\end{equation}
where typically $\omega_E > \omega_C$ for vivid, dread risks.

\textbf{Source:} \citet{loewenstein2000emotions}, \citet{loewenstein2001temptation}
\end{tcolorbox}

\begin{tcolorbox}[colback=gray!5!white, colframe=gray!75!black,
    title=Axiom UNMAPPED_LO-6: Projection Bias]
\textbf{Statement:} People project their current visceral state onto predictions of future preferences.
\begin{equation}
\hat{U}_{t+k}(x | V_t) = \beta \cdot U(x | V_t) + (1-\beta) \cdot U(x | V_{t+k})
\end{equation}
where $\beta > 0$ captures the degree of projection bias.

\textbf{Source:} \citet{loewenstein2000projecting}, \citet{loewenstein2003pain}
\end{tcolorbox}


%%%%%%%%%%%%%%%%%%%%%%%%%%%%%%%%%%%%%%%%%%%%%%%%%%%%%%%%%%%%%%%%%%%%%%%%%%%%%%%
\section{Key Results and Findings}
\label{sec:lo-results}
%%%%%%%%%%%%%%%%%%%%%%%%%%%%%%%%%%%%%%%%%%%%%%%%%%%%%%%%%%%%%%%%%%%%%%%%%%%%%%%

\begin{table}[htbp]
\centering
\caption{Summary of Key Results from Loewenstein's Research}
\label{tab:lo-key-results}
\renewcommand{\arraystretch}{1.3}
\begin{tabular}{@{}llcc@{}}
\toprule
\textbf{Result} & \textbf{Finding} & \textbf{Effect} & \textbf{Tier} \\
\midrule
UNMAPPED_LO-1 & Visceral states override cognition & $\alpha \to 0$ at high V & 1 \\
UNMAPPED_LO-2 & Hot-cold empathy gap exists & 30-50\% underprediction & 1 \\
UNMAPPED_LO-3 & Curiosity peaks at intermediate K & Inverted U-shape & 1 \\
UNMAPPED_LO-4 & Anticipation adds utility & $U_{ant} > 0$ & 1 \\
UNMAPPED_LO-5 & Emotions dominate risk perception & $\omega_E > \omega_C$ & 1 \\
UNMAPPED_LO-6 & Projection bias is systematic & $\beta \approx 0.3-0.5$ & 1 \\
UNMAPPED_LO-7 & Self-serving bias in bargaining & 15-20\% gap & 1 \\
UNMAPPED_LO-8 & Pain memory is reconstructive & Peak-end rule & 1 \\
\bottomrule
\end{tabular}
\end{table}


%%%%%%%%%%%%%%%%%%%%%%%%%%%%%%%%%%%%%%%%%%%%%%%%%%%%%%%%%%%%%%%%%%%%%%%%%%%%%%%
\section{Theme 1: Visceral Influences on Behavior}
\label{sec:lo-visceral}
%%%%%%%%%%%%%%%%%%%%%%%%%%%%%%%%%%%%%%%%%%%%%%%%%%%%%%%%%%%%%%%%%%%%%%%%%%%%%%%

\subsection{The Central Insight}

Loewenstein's foundational paper ``Out of Control'' \citep{loewenstein1996out} establishes that visceral factors---hunger, thirst, sexual arousal, pain, emotions---systematically influence behavior in ways that standard utility theory cannot capture.

\subsection{The Visceral Factor Framework}

Visceral factors share four properties:
\begin{enumerate}[nosep]
\item \textbf{Direct hedonic impact}: They feel good or bad
\item \textbf{Attention capture}: They focus attention on associated goods
\item \textbf{Drive behavior}: They motivate specific actions
\item \textbf{Crowd out cognition}: At high intensity, they override deliberation
\end{enumerate}

\subsection{Mathematical Formalization}

The visceral override function for EBF:
\begin{equation}
A_{\text{eff}}(t^*) = A_{\text{baseline}} \cdot e^{-\lambda V(t^*)}
\end{equation}
where:
\begin{itemize}[nosep]
\item $A_{\text{eff}}$ = Effective awareness (for CORE-AWARE)
\item $V(t^*)$ = Visceral intensity at moment of truth
\item $\lambda$ = Visceral sensitivity parameter ($\approx 2-3$)
\end{itemize}

\subsection{10C Integration}

Visceral factors primarily modify CORE-AWARE and CORE-WHEN:

\begin{table}[htbp]
\centering
\caption{Visceral Factors and 10C Framework}
\label{tab:lo-9c}
\renewcommand{\arraystretch}{1.3}
\begin{tabular}{@{}lll@{}}
\toprule
\textbf{10C Question} & \textbf{Visceral Impact} & \textbf{Mechanism} \\
\midrule
CORE-AWARE ($A$) & Reduces $A$ for long-term utility & Attention capture \\
CORE-WHEN ($\Psi$) & Adds $\Psi_V$ (visceral state) & Context modulation \\
CORE-READY (WAX) & Lowers threshold $\theta$ & Impulse action \\
CORE-STAGE (BCJ) & Can trigger relapse & State-dependent behavior \\
\bottomrule
\end{tabular}
\end{table}


%%%%%%%%%%%%%%%%%%%%%%%%%%%%%%%%%%%%%%%%%%%%%%%%%%%%%%%%%%%%%%%%%%%%%%%%%%%%%%%
\section{Theme 2: Hot-Cold Empathy Gap}
\label{sec:lo-hotcold}
%%%%%%%%%%%%%%%%%%%%%%%%%%%%%%%%%%%%%%%%%%%%%%%%%%%%%%%%%%%%%%%%%%%%%%%%%%%%%%%

\subsection{The Phenomenon}

The hot-cold empathy gap \citep{loewenstein2006hot} describes the systematic failure to predict behavior across visceral states:

\begin{itemize}[nosep]
\item \textbf{Cold-to-hot}: In calm states, people underestimate how much visceral drives will influence them
\item \textbf{Hot-to-cold}: In aroused states, people overestimate how long the state will last
\end{itemize}

\subsection{Applications}

\paragraph{Medical Decision Making.} Patients in current pain choose different treatments than when pain-free. Advance directives made in ``cold'' states may not reflect actual ``hot'' state preferences.

\paragraph{Addiction.} Recovering addicts underestimate craving intensity, leading to relapse. Commitment devices fail because they are designed in ``cold'' states.

\paragraph{Health Behavior.} People planning diet/exercise in comfortable states underestimate the difficulty when hungry or tired.

\subsection{Implications for EBF}

The hot-cold gap implies that parameter estimates may be state-dependent:
\begin{equation}
\theta(\Psi_{\text{hot}}) \neq \theta(\Psi_{\text{cold}})
\end{equation}
This is integrated into CORE-WHERE through state-conditional parameter estimation.


%%%%%%%%%%%%%%%%%%%%%%%%%%%%%%%%%%%%%%%%%%%%%%%%%%%%%%%%%%%%%%%%%%%%%%%%%%%%%%%
\section{Theme 3: Curiosity and Information Gaps}
\label{sec:lo-curiosity}
%%%%%%%%%%%%%%%%%%%%%%%%%%%%%%%%%%%%%%%%%%%%%%%%%%%%%%%%%%%%%%%%%%%%%%%%%%%%%%%

\subsection{The Information Gap Theory}

Loewenstein's curiosity theory \citep{loewenstein1999becausecuriosity} posits that curiosity arises from awareness of an information gap:

\begin{equation}
\text{Curiosity} = f(\text{Gap}) = \text{Gap} \cdot \text{Importance} \cdot \text{Closure Possibility}
\end{equation}

\subsection{Key Predictions}

\begin{enumerate}[nosep]
\item Curiosity is \textit{aversive}---people want to close the gap
\item Curiosity increases with partial knowledge (inverted U)
\item Curiosity is domain-specific, not a general trait
\item Curiosity can override self-interest (e.g., spoilers)
\end{enumerate}

\subsection{EBF Applications}

For behavioral interventions, curiosity can be leveraged:
\begin{itemize}[nosep]
\item \textbf{Engagement}: Create information gaps to maintain attention
\item \textbf{Learning}: Optimal curiosity at intermediate knowledge
\item \textbf{Behavior change}: Frame current state as incomplete
\end{itemize}


%%%%%%%%%%%%%%%%%%%%%%%%%%%%%%%%%%%%%%%%%%%%%%%%%%%%%%%%%%%%%%%%%%%%%%%%%%%%%%%
\section{Theme 4: Anticipatory Utility and Time Preferences}
\label{sec:lo-time}
%%%%%%%%%%%%%%%%%%%%%%%%%%%%%%%%%%%%%%%%%%%%%%%%%%%%%%%%%%%%%%%%%%%%%%%%%%%%%%%

\subsection{Anticipation as Utility Source}

Loewenstein's early work \citep{loewenstein1987emotions} established that anticipation generates utility:

\begin{equation}
U_{\text{total}}(x, T) = \int_0^T \alpha(t) \cdot U_{\text{anticipation}}(x) \, dt + U_{\text{consumption}}(x)
\end{equation}

\subsection{Implications}

\begin{enumerate}[nosep]
\item \textbf{Preference for delay}: Sometimes people prefer to wait (savoring)
\item \textbf{Sequence preferences}: Prefer improving sequences (e.g., save best for last)
\item \textbf{Dread}: Negative anticipation of future pain
\item \textbf{Sign effect}: Different discounting for gains vs. losses
\end{enumerate}

\subsection{Connection to Present Bias}

With Shane Frederick and Ted O'Donoghue, Loewenstein's review \citep{frederick2002time} established the quasi-hyperbolic framework:
\begin{equation}
U_0 = u_0 + \beta \sum_{t=1}^T \delta^t u_t
\end{equation}
where $\beta < 1$ captures present bias.


%%%%%%%%%%%%%%%%%%%%%%%%%%%%%%%%%%%%%%%%%%%%%%%%%%%%%%%%%%%%%%%%%%%%%%%%%%%%%%%
\section{Theme 5: Risk as Feelings}
\label{sec:lo-risk}
%%%%%%%%%%%%%%%%%%%%%%%%%%%%%%%%%%%%%%%%%%%%%%%%%%%%%%%%%%%%%%%%%%%%%%%%%%%%%%%

\subsection{The Risk-as-Feelings Hypothesis}

Loewenstein and colleagues \citep{loewenstein2000emotions} argue that risk perception has two components:

\begin{enumerate}[nosep]
\item \textbf{Cognitive evaluation}: Probability $\times$ outcome
\item \textbf{Emotional reaction}: Fear, dread, anticipatory anxiety
\end{enumerate}

When these diverge, emotional reactions often dominate:
\begin{equation}
\text{Behavior} = \begin{cases}
\text{Cognitive-driven} & \text{if emotion low} \\
\text{Emotion-driven} & \text{if emotion high}
\end{cases}
\end{equation}

\subsection{Predictive Factors for Emotional Dominance}

\begin{itemize}[nosep]
\item \textbf{Vividness}: Concrete, imaginable outcomes trigger emotion
\item \textbf{Immediacy}: Near-term risks feel worse
\item \textbf{Personal relevance}: Affecting self or loved ones
\item \textbf{Dread}: Uncontrollable, catastrophic outcomes
\end{itemize}


%%%%%%%%%%%%%%%%%%%%%%%%%%%%%%%%%%%%%%%%%%%%%%%%%%%%%%%%%%%%%%%%%%%%%%%%%%%%%%%
\section{Integration with Other Appendices}
\label{sec:lo-integration}
%%%%%%%%%%%%%%%%%%%%%%%%%%%%%%%%%%%%%%%%%%%%%%%%%%%%%%%%%%%%%%%%%%%%%%%%%%%%%%%

\subsection{Connection to Appendix ACE (LIT-KAHNEMAN)}

Loewenstein's visceral factors map onto Kahneman's dual systems:

\begin{tcolorbox}[colback=yellow!5!white, colframe=yellow!75!black,
    title=Integration: LIT-LOEWENSTEIN $\leftrightarrow$ LIT-KAHNEMAN]
\textbf{What LIT-LOEWENSTEIN provides:}
\begin{itemize}[nosep]
\item Visceral states as System 1 triggers
\item Hot-cold gap as System 1/System 2 conflict
\item Emotion-driven risk perception
\end{itemize}

\textbf{What LIT-KAHNEMAN provides:}
\begin{itemize}[nosep]
\item Dual-system architecture for processing
\item Heuristics and biases framework
\item Prospect theory reference points
\end{itemize}
\end{tcolorbox}

\subsection{Connection to Appendix UNMAPPED_AU (CORE-AWARE)}

Visceral factors directly modify the Awareness function:
\begin{equation}
A_{\text{Loewenstein}}(\cdot) = A_{\text{baseline}}(\cdot) \cdot g(V) \quad \text{where } g(V) \in [0,1]
\end{equation}

\subsection{Connection to Appendix UNMAPPED_V (CORE-WHEN)}

Visceral states constitute a new $\Psi$ dimension:
\begin{equation}
\Psi = (\Psi_E, \Psi_S, \Psi_T, \Psi_{Sp}, \Psi_I, \Psi_C, \Psi_{Te}, \Psi_{En}, \textbf{$\Psi_V$})
\end{equation}
where $\Psi_V$ captures visceral state intensity.


%%%%%%%%%%%%%%%%%%%%%%%%%%%%%%%%%%%%%%%%%%%%%%%%%%%%%%%%%%%%%%%%%%%%%%%%%%%%%%%
\section{Worked Example: Designing for Visceral States}
\label{sec:lo-example}
%%%%%%%%%%%%%%%%%%%%%%%%%%%%%%%%%%%%%%%%%%%%%%%%%%%%%%%%%%%%%%%%%%%%%%%%%%%%%%%

\subsection{Setup}
A health app wants to help users maintain healthy eating habits. Using Loewenstein's framework to design for visceral state management.

\paragraph{The Problem.} Users commit to healthy eating when not hungry (cold state) but break commitments when hungry (hot state).

\paragraph{Application with Visceral Framework.}

\textbf{Step 1: Identify visceral triggers.}
\begin{itemize}[nosep]
\item Hunger intensity $V_{\text{hunger}} \in [0,1]$
\item Time since last meal
\item Proximity to tempting foods
\end{itemize}

\textbf{Step 2: Design for hot states.}
Rather than relying on cold-state commitments:
\begin{itemize}[nosep]
\item Pre-commitment: Remove unhealthy options before hunger hits
\item Just-in-time intervention: Alert at predicted high-hunger moments
\item Reduce visceral intensity: Suggest healthy snacks to reduce hunger
\end{itemize}

\textbf{Step 3: Address hot-cold empathy gap.}
Help users in cold states anticipate hot states:
\begin{itemize}[nosep]
\item Simulate hunger: ``Imagine it's 6pm and you haven't eaten since noon''
\item Review past failures: ``Last Tuesday at this time, you ate chips''
\item Adjust expectations: ``You will feel strong cravings---this is normal''
\end{itemize}

\paragraph{Result.}
The visceral-aware design addresses:
\begin{itemize}[nosep]
\item Hot-cold gap: Users plan for actual hot-state behavior
\item Visceral override: Interventions reduce $V$ before it gets too high
\item Projection bias: Explicit reminders of state-dependent preferences
\end{itemize}

\paragraph{Discussion.}
This example shows how EBF can use Loewenstein's framework to design interventions that account for visceral influences, rather than assuming purely cognitive decision-making.


%%%%%%%%%%%%%%%%%%%%%%%%%%%%%%%%%%%%%%%%%%%%%%%%%%%%%%%%%%%%%%%%%%%%%%%%%%%%%%%
\section{Implications for Practice}
\label{sec:lo-practice}
%%%%%%%%%%%%%%%%%%%%%%%%%%%%%%%%%%%%%%%%%%%%%%%%%%%%%%%%%%%%%%%%%%%%%%%%%%%%%%%

\begin{tcolorbox}[colback=green!5!white, colframe=green!75!black,
    title=Practical Implications]
\begin{enumerate}
\item \textbf{Design for hot states:} Don't rely on cold-state commitments; create environments that work when willpower is depleted
\item \textbf{Bridge the empathy gap:} Help people in cold states accurately predict hot-state behavior through simulation and past data
\item \textbf{Leverage curiosity:} Create information gaps to drive engagement; close gaps to drive satisfaction
\item \textbf{Use anticipation:} Build anticipatory utility into experiences; sequence events for improving trajectories
\item \textbf{Address emotional risk:} Communicate risks in ways that align cognitive and emotional responses
\end{enumerate}
\end{tcolorbox}


%%%%%%%%%%%%%%%%%%%%%%%%%%%%%%%%%%%%%%%%%%%%%%%%%%%%%%%%%%%%%%%%%%%%%%%%%%%%%%%
\section{Parameters for LLMMC Estimation}
\label{sec:lo-llmmc}
%%%%%%%%%%%%%%%%%%%%%%%%%%%%%%%%%%%%%%%%%%%%%%%%%%%%%%%%%%%%%%%%%%%%%%%%%%%%%%%

From Loewenstein's research, the following parameters are available for LLMMC:

\begin{table}[htbp]
\centering
\caption{LLMMC Parameters from LIT-LOEWENSTEIN}
\label{tab:lo-llmmc}
\renewcommand{\arraystretch}{1.3}
\begin{tabular}{@{}llccc@{}}
\toprule
\textbf{Parameter} & \textbf{Description} & \textbf{Point Estimate} & \textbf{95\% CI} & \textbf{Tier} \\
\midrule
$\lambda_V$ & Visceral sensitivity & 2.5 & [2.0, 3.0] & 1 \\
$\beta_{\text{proj}}$ & Projection bias & 0.4 & [0.3, 0.5] & 1 \\
Gap$_{\text{HC}}$ & Hot-cold empathy gap & 0.35 & [0.30, 0.50] & 1 \\
$\omega_E / \omega_C$ & Emotion/cognition ratio & 1.5 & [1.2, 2.0] & 1 \\
$\beta$ & Present bias (quasi-hyp.) & 0.7 & [0.6, 0.8] & 1 \\
\bottomrule
\end{tabular}
\end{table}


%%%%%%%%%%%%%%%%%%%%%%%%%%%%%%%%%%%%%%%%%%%%%%%%%%%%%%%%%%%%%%%%%%%%%%%%%%%%%%%
\section{Summary}
\label{sec:lo-summary}
%%%%%%%%%%%%%%%%%%%%%%%%%%%%%%%%%%%%%%%%%%%%%%%%%%%%%%%%%%%%%%%%%%%%%%%%%%%%%%%

\begin{tcolorbox}[colback=green!10!white, colframe=green!75!black,
    title=\textbf{Appendix UNMAPPED_LO: Summary}]

\textbf{Core Question:} How do visceral states and emotions override cognitive preferences?

\textbf{Main Contributions:}
\begin{itemize}[nosep]
    \item Visceral override: Hot states dominate cold preferences
    \item Hot-cold empathy gap: Systematic prediction failures (30-50\%)
    \item Curiosity as information gap: Inverted U-shape
    \item Anticipatory utility: Value from expecting future events
    \item Risk as feelings: Emotions drive risk perception
    \item Projection bias: Current state projected onto future ($\beta \approx 0.4$)
\end{itemize}

\textbf{Key Outputs:}
\begin{itemize}[nosep]
    \item $V$: Visceral intensity parameter
    \item $\lambda_V$: Visceral sensitivity ($\approx 2.5$)
    \item Gap$_{\text{HC}}$: Hot-cold empathy gap ($\approx 35\%$)
\end{itemize}

\textbf{Integration:} This appendix connects to CORE-AWARE (visceral modulation of $A$),
CORE-WHEN ($\Psi_V$ dimension), and LIT-KAHNEMAN via System 1/System 2 architecture.

\end{tcolorbox}


%%%%%%%%%%%%%%%%%%%%%%%%%%%%%%%%%%%%%%%%%%%%%%%%%%%%%%%%%%%%%%%%%%%%%%%%%%%%%%%
\section{Glossary of Symbols}
\label{sec:lo-glossary}
%%%%%%%%%%%%%%%%%%%%%%%%%%%%%%%%%%%%%%%%%%%%%%%%%%%%%%%%%%%%%%%%%%%%%%%%%%%%%%%

\begin{tcolorbox}[colback=blue!5!white, colframe=blue!75!black]
\textbf{Central Reference:} For complete symbol definitions, see
\textbf{Appendix G} (\textit{Glossary and Notation}).

This section lists only symbols introduced or specialized in this appendix.
\end{tcolorbox}

\vspace{0.5em}

\begin{table}[htbp]
\centering
\caption{Symbols Introduced in Appendix UNMAPPED_LO}
\label{tab:lo-symbols}
\renewcommand{\arraystretch}{1.3}
\begin{tabular}{@{}clp{6cm}c@{}}
\toprule
\textbf{Symbol} & \textbf{Name} & \textbf{Definition} & \textbf{Also in G?} \\
\midrule
$V$ & Visceral intensity & Intensity of visceral state $\in [0,1]$ & NEW \\
$\lambda_V$ & Visceral sensitivity & How quickly visceral states override cognition & NEW \\
$\Psi_V$ & Visceral context & Visceral state as context dimension & NEW \\
Gap$_{\text{HC}}$ & Hot-cold gap & Empathy gap magnitude & NEW \\
$\beta_{\text{proj}}$ & Projection bias & Degree of current-state projection & NEW \\
$\omega_E, \omega_C$ & Emotion/cognition weights & Weights in risk perception & NEW \\
$A$ & Awareness function & (Modified by visceral state) & \checkmark \\
\bottomrule
\end{tabular}
\end{table}


%%%%%%%%%%%%%%%%%%%%%%%%%%%%%%%%%%%%%%%%%%%%%%%%%%%%%%%%%%%%%%%%%%%%%%%%%%%%%%%
\section{Critical Foundations and Objections}
\label{sec:lo-foundations}
%%%%%%%%%%%%%%%%%%%%%%%%%%%%%%%%%%%%%%%%%%%%%%%%%%%%%%%%%%%%%%%%%%%%%%%%%%%%%%%

\subsection{Objection 1: Visceral Effects Are Just Noise}

\textbf{Objection:} Visceral influences are random perturbations, not systematic biases. Standard utility theory can accommodate them as preference shocks.

\textbf{Response:} Loewenstein shows visceral effects are:
\begin{itemize}[nosep]
\item \textit{Systematic}: Predictable direction (toward immediate gratification)
\item \textit{State-dependent}: Consistently tied to measurable physiological states
\item \textit{Cross-situational}: Same pattern across hunger, pain, arousal, etc.
\end{itemize}

\subsection{Objection 2: Laboratory Artifacts}

\textbf{Objection:} Hot-cold effects are laboratory artifacts that don't predict field behavior.

\textbf{Response:} Field evidence supports the theory:
\begin{itemize}[nosep]
\item Medical decisions vary with current pain \citep{loewenstein2006hot}
\item Grocery shopping varies with hunger (Nisbett \& Kanouse 1969)
\item Addiction relapse follows visceral craving patterns
\end{itemize}

\subsection{Objection 3: Tautological Theory}

\textbf{Objection:} ``Visceral states affect behavior'' is tautological---of course bodily states matter.

\textbf{Response:} The contribution is:
\begin{itemize}[nosep]
\item \textit{Specific predictions}: Hot-cold gap magnitude, inverted-U curiosity
\item \textit{Design implications}: How to build commitment devices that work
\item \textit{Formalization}: Mathematical specification enabling integration with EBF
\end{itemize}


%%%%%%%%%%%%%%%%%%%%%%%%%%%%%%%%%%%%%%%%%%%%%%%%%%%%%%%%%%%%%%%%%%%%%%%%%%%%%%%
\section{Open Issues and Research Agenda}
\label{sec:lo-open}
%%%%%%%%%%%%%%%%%%%%%%%%%%%%%%%%%%%%%%%%%%%%%%%%%%%%%%%%%%%%%%%%%%%%%%%%%%%%%%%

\begin{table}[htbp]
\centering
\caption{Research Roadmap for Loewenstein-Related Questions}
\label{tab:lo-roadmap}
\renewcommand{\arraystretch}{1.3}
\begin{tabular}{@{}clcl@{}}
\toprule
\textbf{\#} & \textbf{Issue} & \textbf{Priority} & \textbf{Status} \\
\midrule
1 & Individual differences in visceral sensitivity & HIGH & Partial evidence \\
2 & Cultural variation in hot-cold gap & MEDIUM & Limited data \\
3 & Neurobiological mechanisms of override & HIGH & Active research \\
4 & Long-term adaptation to visceral awareness & MEDIUM & Underexplored \\
5 & Interaction between visceral factors & LOW & Speculative \\
\bottomrule
\end{tabular}
\end{table}


%%%%%%%%%%%%%%%%%%%%%%%%%%%%%%%%%%%%%%%%%%%%%%%%%%%%%%%%%%%%%%%%%%%%%%%%%%%%%%%
\section{Complete Paper Inventory}
\label{sec:lo-papers}
%%%%%%%%%%%%%%%%%%%%%%%%%%%%%%%%%%%%%%%%%%%%%%%%%%%%%%%%%%%%%%%%%%%%%%%%%%%%%%%

\begin{table}[htbp]
\centering
\caption{All 29 Loewenstein Papers in EBF Database}
\label{tab:lo-papers}
\renewcommand{\arraystretch}{1.2}
\scriptsize
\begin{tabular}{@{}llll@{}}
\toprule
\textbf{ID} & \textbf{Year} & \textbf{Theme} & \textbf{Tier} \\
\midrule
loewenstein1987emotions & 1987 & Anticipatory Utility & 1 \\
prelec1991puzzles & 1991 & Time Preferences & 1 \\
PAP-loewenstein1996out & 1992 & Visceral Influences & 1 \\
babcock1995wages & 1995 & Self-Serving Bias & 1 \\
loewenstein1996out & 1996 & Visceral Influences & 1 \\
loewenstein1998anticipatory & 1998 & Emotions & 1 \\
loewenstein1999becausecuriosity & 1999 & Curiosity & 1 \\
loewenstein2000emotions & 2000 & Risk as Feelings & 2 \\
loewenstein2000projecting & 2000 & Projection Bias & 1 \\
loewenstein2000shame & 2000 & Emotions & 1 \\
ariely2001cheating & 2001 & (with Ariely) & 1 \\
loewenstein2001temptation & 2001 & Temptation & 1 \\
frederick2002time & 2002 & Time Preferences & 3 \\
loewenstein2003pain & 2003 & Pain Decisions & 1 \\
ariely2004coherentanchoring & 2004 & (with Ariely) & 1 \\
PAP-camerer2004neuroeconomicsneuroeconomics & 2004 & (with Camerer) & 1 \\
loewenstein2005visceral & 2005 & Visceral Influences & 1 \\
loewenstein2006hot & 2006 & Hot-Cold Gap & 1 \\
ariely2007expensiveexpensive & 2007 & (with Ariely) & 1 \\
loewenstein2007should & 2007 & Behavioral Models & 1 \\
loewenstein2008choice & 2008 & Emotions & 1 \\
loewenstein2009present & 2009 & Present Bias & 1 \\
loewenstein2010visceral & 2010 & Neuroeconomics & 1 \\
camerer2011trust & 2011 & (with Camerer) & 1 \\
loewenstein2012progress & 2012 & Behavioral Econ & 1 \\
ariely2013the & 2013 & (with Ariely) & 1 \\
loewenstein2013visceral & 2013 & Visceral States & 1 \\
loewenstein2014visceral & 2014 & Self-Control & 1 \\
ariely2016pain & 2016 & (with Ariely) & 1 \\
\bottomrule
\end{tabular}
\end{table}


%%%%%%%%%%%%%%%%%%%%%%%%%%%%%%%%%%%%%%%%%%%%%%%%%%%%%%%%%%%%%%%%%%%%%%%%%%%%%%%
\section{References}
\label{sec:lo-references}
%%%%%%%%%%%%%%%%%%%%%%%%%%%%%%%%%%%%%%%%%%%%%%%%%%%%%%%%%%%%%%%%%%%%%%%%%%%%%%%

\begin{tcolorbox}[colback=blue!5!white, colframe=blue!75!black]
\textbf{Master Bibliography:} For complete reference details, see
\texttt{bcm\_master.bib} in the repository.

This section lists key references for this appendix.
\end{tcolorbox}

\vspace{0.5em}

\subsection{Key References}

\begin{itemize}[nosep]
    \item \textbf{Visceral Influences:} \citet{loewenstein1996out}, \citet{loewenstein2005visceral}
    \item \textbf{Hot-Cold Gap:} \citet{loewenstein2006hot}, \citet{loewenstein2000projecting}
    \item \textbf{Curiosity:} \citet{loewenstein1999becausecuriosity}
    \item \textbf{Anticipatory Utility:} \citet{loewenstein1987emotions}
    \item \textbf{Risk as Feelings:} \citet{loewenstein2000emotions}
    \item \textbf{Time Preferences:} \citet{frederick2002time}
\end{itemize}

% Link to master bibliography
\nocite{bcm_master}


%%%%%%%%%%%%%%%%%%%%%%%%%%%%%%%%%%%%%%%%%%%%%%%%%%%%%%%%%%%%%%%%%%%%%%%%%%%%%%%
% CROSS-REFERENCE FOOTER (Required)
%%%%%%%%%%%%%%%%%%%%%%%%%%%%%%%%%%%%%%%%%%%%%%%%%%%%%%%%%%%%%%%%%%%%%%%%%%%%%%%

\vspace{1em}
\hrule
\vspace{0.5em}

\noindent\textit{Cross-references:}
Appendix G (Central Glossary),
Appendix ACE (LIT-KAHNEMAN),
Appendix UNMAPPED_AU (CORE-AWARE),
Appendix UNMAPPED_V (CORE-WHEN),
Appendix UNMAPPED_BBB (CORE-WHERE).

\noindent\textit{Master Bibliography:} \texttt{bcm\_master.bib}


% =============================================================================
% END OF APPENDIX LO
% =============================================================================
